% GENERATED FILE. Edit canonical lessons, then run npm run generate:books.
% canonical-source-hash: fnv1a64:61fea796a25b1d71
% canonical-lessons: GU-C111-kurto, GU-C111-payjamo, GU-C111-viti, GU-C111-butti, GU-C111-moja

\chapter{Kurta, Pyjamas, Ring, Earring, Socks}
\label{ch:gu-kurta-pyjamas-ring-earring-socks}
\hlchaptermodality{111}

\begin{chapteropening}
\textbf{By the end of this chapter:} \emph{I can say a kurta, pyjamas, a ring, an earring and socks.}

Complete the last lesson of chapter 111: I can say a kurta, pyjamas, a ring, an earring and socks.
\end{chapteropening}

\section[\texorpdfstring{\gu{કુરતો} (kurto)}{કુરતો (kurto)}]{\gu{કુરતો} (kurto) — a kurta}

\label{lesson:GU-C111-kurto}



\begin{quote}
Before the new one: say the Gujarati for a pillow, then the Gujarati for a doll.
\end{quote}

\subsection*{You'll want to know: \gu{કુરતો}}
\textbf{\gu{કુરતો}} (\emph{kurto}) — \textquotedblleft{}a kurta\textquotedblright{}.

\begin{cousinweb}
\textbf{\gu{કુરતો}} (\emph{kurto}) is a loan, usually traced to Persian \emph{kurta}, \textquotedblleft{}a long shirt\textquotedblright{}.
\end{cousinweb}

\begin{grammarlens}[title={What you've built}]
One more word for this chapter.
\end{grammarlens}

\subsection*{Your turn}
\begin{itemize}
\raggedright
  \item \emph{Say it:} \emph{kurto}
  \item \emph{Say it:} \emph{kurto}, once more
  \item \emph{Say it:} say \emph{kurto}
  \item \emph{Recall:} say the Gujarati for a pillow, then the Gujarati for a doll, then say \emph{kurto} again
\end{itemize}

\begin{tcolorbox}[breakable,colback=teal!4,colframe=teal!35!black,title={Before you move on}]
What does \gu{કુરતો} mean? (\textquotedblleft{}A kurta\textquotedblright{}.) Say it once more.
\end{tcolorbox}
\section[\texorpdfstring{\gu{પાયજામો} (pāyjāmo)}{પાયજામો (pāyjāmo)}]{\gu{પાયજામો} (pāyjāmo) — pyjamas}

\label{lesson:GU-C111-payjamo}



\begin{quote}
Before the new one: say the Gujarati for a doll, then the Gujarati for a kurta.
\end{quote}

\subsection*{You'll want to know: \gu{પાયજામો}}
\textbf{\gu{પાયજામો}} (\emph{pāyjāmo}) — \textquotedblleft{}pyjamas\textquotedblright{}.

\begin{cousinweb}
\textbf{\gu{પાયજામો}} (\emph{pāyjāmo}) is a loan, usually traced to Persian \emph{pāy-jāma}, \textquotedblleft{}leg garment\textquotedblright{}: \emph{pāy}, \textquotedblleft{}leg\textquotedblright{}, and \emph{jāma}, \textquotedblleft{}garment\textquotedblright{}.
\end{cousinweb}

\begin{grammarlens}[title={What you've built}]
One more word for this chapter.
\end{grammarlens}

\subsection*{Your turn}
\begin{itemize}
\raggedright
  \item \emph{Say it:} \emph{pāyjāmo}
  \item \emph{Say it:} \emph{pāyjāmo}, once more
  \item \emph{Say it:} say \emph{pāyjāmo}
  \item \emph{Recall:} say the Gujarati for a doll, then the Gujarati for a kurta, then say \emph{pāyjāmo} again
\end{itemize}

\begin{tcolorbox}[breakable,colback=teal!4,colframe=teal!35!black,title={Before you move on}]
What does \gu{પાયજામો} mean? (\textquotedblleft{}Pyjamas\textquotedblright{}.) Say it once more.
\end{tcolorbox}
\section[\texorpdfstring{\gu{વીંટી} (vī̃ṭī)}{વીંટી (vī̃ṭī)}]{\gu{વીંટી} (vī̃ṭī) — a ring}

\label{lesson:GU-C111-viti}



\begin{quote}
Before the new one: say the Gujarati for a kurta, then the Gujarati for pyjamas.
\end{quote}

\subsection*{You'll want to know: \gu{વીંટી}}
\textbf{\gu{વીંટી}} (\emph{vī̃ṭī}) — \textquotedblleft{}a ring\textquotedblright{}.

\begin{cousinweb}
CDIAL lists \textbf{\gu{વીંટી}} (\emph{vīṁṭī}), \textquotedblleft{}ring\textquotedblright{}, under Sanskrit \emph{vīṭā}, \textquotedblleft{}small piece of stick; roll of betel\textquotedblright{}, heading a loose group of words alike in shape and sense (CDIAL 12045). Austroasiatic links have been proposed for some of them.
\end{cousinweb}

\begin{grammarlens}[title={What you've built}]
One more word for this chapter.
\end{grammarlens}

\subsection*{Your turn}
\begin{itemize}
\raggedright
  \item \emph{Say it:} \emph{vī̃ṭī}
  \item \emph{Say it:} \emph{vī̃ṭī}, once more
  \item \emph{Say it:} say \emph{vī̃ṭī}
  \item \emph{Recall:} say the Gujarati for a kurta, then the Gujarati for pyjamas, then say \emph{vī̃ṭī} again
\end{itemize}

\begin{tcolorbox}[breakable,colback=teal!4,colframe=teal!35!black,title={Before you move on}]
What does \gu{વીંટી} mean? (\textquotedblleft{}A ring\textquotedblright{}.) Say it once more.
\end{tcolorbox}
\section[\texorpdfstring{\gu{બુટ્ટી} (buṭṭī)}{બુટ્ટી (buṭṭī)}]{\gu{બુટ્ટી} (buṭṭī) — an earring}

\label{lesson:GU-C111-butti}



\begin{quote}
Before the new one: say the Gujarati for pyjamas, then the Gujarati for a ring.
\end{quote}

\subsection*{You'll want to know: \gu{બુટ્ટી}}
\textbf{\gu{બુટ્ટી}} (\emph{buṭṭī}) — \textquotedblleft{}an earring\textquotedblright{}.

\begin{cousinweb}
The origin of \textbf{\gu{બુટ્ટી}} (\emph{buṭṭī}) is not settled. CDIAL lists a Gujarati \emph{buṭṭī} only as \textquotedblleft{}painted flower\textquotedblright{}, under a reconstructed \emph{*būṭa}, \textquotedblleft{}bush, plant\textquotedblright{} (CDIAL 9297); whether the earring is the same word is not known.
\end{cousinweb}

\begin{grammarlens}[title={What you've built}]
One more word for this chapter.
\end{grammarlens}

\subsection*{Your turn}
\begin{itemize}
\raggedright
  \item \emph{Say it:} \emph{buṭṭī}
  \item \emph{Say it:} \emph{buṭṭī}, once more
  \item \emph{Say it:} say \emph{buṭṭī}
  \item \emph{Recall:} say the Gujarati for pyjamas, then the Gujarati for a ring, then say \emph{buṭṭī} again
\end{itemize}

\begin{tcolorbox}[breakable,colback=teal!4,colframe=teal!35!black,title={Before you move on}]
What does \gu{બુટ્ટી} mean? (\textquotedblleft{}An earring\textquotedblright{}.) Say it once more.
\end{tcolorbox}
\section[\texorpdfstring{\gu{મોજાં} (mojā̃)}{મોજાં (mojā̃)}]{\gu{મોજાં} (mojā̃) — socks}

\label{lesson:GU-C111-moja}



\begin{quote}
Before the new one: say the Gujarati for a ring, then the Gujarati for an earring.
\end{quote}

\subsection*{You'll want to know: \gu{મોજાં}}
\textbf{\gu{મોજાં}} (\emph{mojā̃}) — \textquotedblleft{}socks\textquotedblright{}.

\begin{cousinweb}
\textbf{\gu{મોજાં}} (\emph{mojāṁ}) is a loan, usually traced to Persian \emph{mōza}, \textquotedblleft{}sock, stocking\textquotedblright{}.
\end{cousinweb}

\begin{grammarlens}[title={What you've built}]
One more word for this chapter.
\end{grammarlens}

\subsection*{Your turn}
\begin{itemize}
\raggedright
  \item \emph{Say it:} \emph{mojā̃}
  \item \emph{Say it:} \emph{mojā̃}, once more
  \item \emph{Say it:} say \emph{mojā̃}
  \item \emph{Recall:} say the Gujarati for a ring, then the Gujarati for an earring, then say \emph{mojā̃} again
\end{itemize}

\begin{tcolorbox}[breakable,colback=teal!4,colframe=teal!35!black,title={Before you move on}]
What does \gu{મોજાં} mean? (\textquotedblleft{}Socks\textquotedblright{}.) Say it once more.
\end{tcolorbox}
